\section{Validation and error handling}
\label{sec:errors}

Three points are for ``comprehensive handling and validation across front and back end; all inputs
validated; errors handled gracefully with clear messages; prevents invalid data''.

\subsection{Principles}

\begin{enumerate}
  \item \textbf{Both tiers check, from one table.} The browser checks so that the user is told at
    once; the server checks again because a request need not come from our browser.
  \item \textbf{Never invent a result.} An input that cannot be scored gets an explanation, not a
    score (AC-02).
  \item \textbf{Every state has a next action}, and no error clears what the user typed.
  \item \textbf{Messages are specific and in plain words}: what is wrong, and what to do.
  \item \textbf{One body, one place.} The server builds every error in its handlers; the client
    turns every failure into one \code{ApiError}; views switch on its \code{status}.
\end{enumerate}

\subsection{The validation table}
\label{sec:validation}

{\small
\begin{longtable}{@{}L{2.2cm}L{5.5cm}L{3.2cm}L{3.2cm}@{}}
\toprule
\textbf{Input} & \textbf{Rule} & \textbf{Front end} & \textbf{Back end} \\
\midrule
\endhead
Bitcoin address & Not empty; at most 100 characters; no spaces. Not an Ethereum or test-network address (each has its own message). Legacy: 26 to 35 characters of the Base58 alphabet, a correct checksum, version 0 or 5. \code{bc1}: one letter case, the Bech32 alphabet, a correct Bech32 or Bech32m checksum, a valid program length & \fp{lib/address.ts} (exists), as the user types and on submit & \fp{core/validation.py}, on E4 and on every path with an address; 400 \\
Day & \code{YYYY-MM-DD} and a real date & Date fields limited to the address's active period & Parsed; 400 \\
Day range & \code{from} not after \code{to}; not together with \code{date} & Both fields marked & 400 with both fields named \\
Transaction ID & 64 hexadecimal characters & Checked when a link is opened & Pattern; 400. Unknown for that address: 404 \\
Amount & A number; greater than 0; at least 1 satoshi; at most 21,000,000 BTC; at most 8 decimal places in BTC, whole in satoshis. One comma is read as the decimal point; grouping separators are refused with a message & As the user types & Parsed as a decimal, never a float; 400 with \code{fields} \\
Direction & \code{in} or \code{out} & A two-way switch; cannot be empty & Choice; 400 \\
Time & A real date and time in UTC; not before the address's first transaction; at most \code{maxGapYears} after its last & Fields limited to the allowed range & 422 if too early to score; 400 if too far ahead \\
List parameters & \code{page} from 1; \code{pageSize} 1 to 200; \code{sort}, \code{status}, \code{direction} and \code{method} from their lists & Built by the client, never typed & 400 naming the parameter \\
Watchlist label and note & 40 and 500 characters at most; no control characters; kept and shown as plain text & Character counters; Save disabled when over & 400; SQL through parameters only \\
Export format & \code{csv} or \code{json} & A menu & Choice; 400 \\
Live lookup & A valid address not in the snapshot; \code{maxTx} 50 to 1,000 & The dialog & 400; 409; 429 \\
Any body & JSON; at most 10 kB; no unknown fields & -- & 415; 413; 400 \\
\bottomrule
\end{longtable}
}

\textbf{Shared test vectors.} \fp{shared/address_vectors.json} holds valid addresses of each type
(P2PKH, P2SH, P2WPKH, P2WSH, P2TR) and one invalid address for each failure the checker names.
Vitest and pytest both read it, so the two tiers must agree on every verdict and message key. The
same is done for the what-if rules with \fp{shared/whatif_vectors.json}.

\subsection{The catalogue of states}
\label{sec:catalogue}

The first six are Assignment 1's (FR-05 and the network state). Each row has a component test and,
where a server status is involved, an API test.

{\small
\begin{longtable}{@{}L{2.9cm}L{2.7cm}L{4.9cm}L{3.6cm}@{}}
\toprule
\textbf{Situation} & \textbf{Signal} & \textbf{What the user sees} & \textbf{Next action} \\
\midrule
\endhead
Invalid input & Client rule, or 400 & The specific message under the field; the field keeps its text and focus & Correct it \\
Unsupported address & 404 \code{unsupported} & What the snapshot covers; three supported examples & Open an example; clear the box; with S-09, fetch it \\
Loading & Request pending & A skeleton in the final layout; ``Still working'' after 5 s & Wait; leave \\
Empty & 200, no rows & ``Choose a transaction to see its details''; ``No transaction matches these filters'' & Choose one; reset the filters \\
Insufficient history & 422, or \code{history.status} & The transactions unscored, and why: ``8 transactions in under a day; scoring needs 10 over 14 days'' & Browse them anyway; choose another address; in V4, set the first date that works \\
Network failure & No answer & ``The explorer service did not respond.'' The status dot turns red & Try again; reloads by itself when the service is back \\
Not found & 404 \code{not_found}, or an unknown route & ``There's nothing at this link'' & Go to the address, or to the list \\
Timeout & 10 s (90 s for the Assistant) & ``This is taking longer than it should'' & Try again \\
Server error & 500 & ``Something went wrong on our side'', with the request identifier & Try again; copy the identifier \\
Service starting & 503 & ``The service is starting; this page will load by itself'' & None needed \\
Busy & 429 & The wait, counting down & Retries by itself \\
Upstream failure (S-09) & 502, 504 & ``The public blockchain service did not answer'' & Try again; pick a snapshot address \\
Conflict & 409 & ``The watchlist is full (200)''; ``This address is already being fetched'' & Remove an entry; open the running lookup \\
Chart failed to draw & Caught by the chart's boundary & The table view, and a note & Show the chart again \\
Part of a view crashed & Error boundary & ``This part of the page failed to load''; the rest keeps working & Reload the part \\
Offline & The browser reports it & A banner; views already loaded stay usable & Reconnect \\
Assistant not installed, or its answer rejected & \code{skippedWriters}, \code{attempts} & A note; the template explanation is shown & Read the README; nothing else \\
\bottomrule
\end{longtable}
}

\subsection{On the server}

\begin{itemize}
  \item \textbf{Handlers.} One each for request validation (400), the package's
    \code{Unsupported} (404) and \code{InsufficientHistory} (422), \code{ValueError} (400),
    missing files (503), HTTP errors raised by routes, and everything else (500).
  \item \textbf{Logging.} One line per request: identifier, method, path, status, milliseconds.
    A 500 logs its trace under the same identifier. Bodies are not logged.
  \item \textbf{What a response never contains.} A stack trace, a file path, a library's raw
    error text, or a banned word.
  \item \textbf{Start-up.} A missing model file, a failed hash or a wrong library version stops the
    start with one sentence that says what to do; it does not start half-working.
\end{itemize}

\subsection{Hardening}

\begin{itemize}
  \item User text never reaches a file path, a shell or an SQL string: addresses are looked up
    in memory, cache files are named by a hash, and SQL uses parameters.
  \item The server calls out only to the two Esplora hosts named in the settings, and only with an
    address that has passed validation.
  \item CORS allows only the listed origins. Bodies over 10 kB are refused. Live lookups and the
    Assistant are rate-limited.
  \item React renders every string as text; labels and notes are never treated as markup. External
    links open with \code{rel="noopener noreferrer"}.
  \item CSV cells that could run as formulas are neutralised (E10).
  \item Dependencies are pinned; \code{npm audit} and \code{pip-audit} run before the code freeze.
\end{itemize}

\begin{done}
Every row of the validation table has a test on each tier and every row of the catalogue has a
component test. Posting each invalid vector to the API returns its expected status and never a 500.
The video shows four error states and one recovery.
\end{done}
